% Original: PDF strana 3, donji slajd.
\slajdid{02-s006}
\begin{frame}[label=02-s006]{Kada je izlaz kombinacione mreže važeći?}
% element: 02-s006:e01
\par\Rezultat{Izlaz zavisi od „trenutnih“ ulaza tek \alert{po završetku prelaznih procesa}.}
\par\vspace{5mm}
% element: 02-s006:e02
\begin{enumerate}
\item Promena ulaza se prenosi kroz pojedina logička kola sa kašnjenjem.
\item Vreme do važećeg izlaza zavisi od kašnjenja duž putanja kroz mrežu.
\end{enumerate}
\par\vspace{5mm}
% element: 02-s006:e03
\Oznaka{Neželjeni efekat}\par\vspace{2mm}
U prelaznom intervalu mogu se pojaviti \alert{lažne nule i jedinice} na izlazu.\par\vspace{3mm}
Ovoj pojavi posvetićemo posebnu pažnju pri sintezi kombinacionih mreža.
\end{frame}
